(1) 球心速度 \(a\dot\theta\)，无滑移使球相对质心自旋速率大小为 \(a\dot\theta/r\)，总动能为 \(\frac12m(1+\beta)a^2\dot\theta^2\)。重力势能取最低点为零 \(U=mga(1-\cos\theta)\)，Euler–Lagrange方程因此是 \((1+\beta)a\ddot\theta=-g\sin\theta\)。
(2) 在 \(|\theta|\ll1\) 时 \(\sin\theta\simeq\theta\)，故小振动频率 \(\omega_0=\sqrt{g/[(1+\beta)a]}\)，周期 \(T=2\pi\sqrt{(1+\beta)(R-r)/g}\)；\(\beta\) 越大意味着同一平移速度下额外转动动能越多，振动越慢。
(3) 沿指向轨道圆心的法向应用牛顿定律，壁面对球的法向支持力为 \(N=m[a\dot\theta^2+g\cos\theta]\)。由能量 \(\frac12m(1+\beta)a^2\dot\theta^2+mga(1-\cos\theta)=\frac12m(1+\beta)v_0^2\)，得到 \(N/m=v_0^2/a\, -[2g/(1+\beta)](1-\cos\theta)+g\cos\theta\)；沿运动轨迹若首次出现 \(N=0\) 且随后需要负支持力，则不能保持原几何接触。
(4) 在最低点静止时 \(v_0=0,\theta=0\)，\(N=mg\)；小振幅始终有 \(N>0\)。形式上 \(\beta\to0\) 得无转动惯量的滑珠角频率 \(\sqrt{g/a}\)。无滑移约束另需足够静摩擦以提供转矩，而 \(N\ge0\) 仅保证法向可接触，两者须分别检查。
